\documentclass[twoside]{article}

% morewrites virtualises \newwrite so a document that combines many glossary
% groups (each costs one \write), biblatex (\bcfout), hyperref (\@outlinefile),
% the toc, and the LOT/LOF can exceed XeLaTeX's 16-register limit. Must be
% loaded before any package that allocates writes; that is why it sits ahead
% of xcolor here.
\usepackage{morewrites}

\usepackage[svgnames,dvipsnames,x11names]{xcolor}
\usepackage{graphicx}
% Babel must be loaded before any package that hooks into its localisation (notably
% the glossaries fragment in ``extra_packages``), otherwise the default acronym
% table title falls back to English regardless of ``language: french`` etc.
\usepackage[english]{babel}
% Hyperref must be loaded before ``glossaries`` (pulled in via ``extra_packages``)
% so glossaries can wrap glossary page numbers in ``\hyperpage`` for clickable
% backreferences. Loading it after glossaries leaves the page numbers as plain
% text without hyperlinks.
\usepackage[colorlinks=true, linkcolor=NavyBlue, urlcolor=NavyBlue, citecolor=NavyBlue]{hyperref}
\usepackage[margin=15mm]{geometry}
\usepackage{ts-fonts}
\usepackage{float}
\usepackage{seqsplit}
\usepackage{ts-callouts}

\usepackage{lastpage} % pour le label LastPage
\usepackage{refcount} % pour \getpagerefnumber
\makeatletter
\def\ps@plain{%
  \let\@oddhead\@empty
  \let\@evenhead\@empty
  \def\@oddfoot{%
    \hfil
    % Afficher le numéro seulement si le document a plus d'une page
    \ifnum\getpagerefnumber{LastPage}>1\relax
      \thepage
    \fi
    \hfil}%
  \let\@evenfoot\@oddfoot
}
\makeatother

\title{Admonitions / Callouts\\[0.5em]\large An Overview of \textbf{fancy} framed elements}
\author{}\date{}
% In two-column mode an overfull line bleeds into the adjacent column, which is
% far worse than a slightly loose line.  \emergencystretch lets TeX add extra
% inter-word stretch as a last resort before declaring an overfull box.
\setlength{\emergencystretch}{3em}
% Forbid widows and orphans of 1 or 2 lines pushed to the next column/page.
% \widowpenalties N p1..pN: p1=penalty when last 1 line is alone at top,
% p2=when last 2 lines are alone, etc.  10000 = forbidden, 150 = default.
% \clubpenalties: same logic for first lines left alone at the bottom.
\widowpenalties 3 10000 10000 150
\clubpenalties  3 10000 10000 150
\displaywidowpenalty 10000
% Restore the single-column abstract style (small font + quotation indentation).
% The article class degrades \begin{abstract} to a bare \section* when the
% twocolumn class option is in effect; we manage column switching manually
% below with \twocolumn[...] so that option is not set here.
\renewenvironment{abstract}{%
  \small
  \begin{center}{\bfseries\abstractname}\end{center}%
  \begin{quotation}%
}{\end{quotation}}
% \thanks{} footnote fix for two-column layout.
% Three bugs defeat the naive toks approach:
%   1. \the\c@footnote is stored literally in toks, not expanded — at replay
%      time article.cls has reset c@footnote to 0 via \@addtoreset{footnote}{page}
%      triggered by the \shipout inside \twocolumn.
%   2. \noexpand before \footnotetext in a toks list corrupts \@ifnextchar
%      scanning when the list is replayed.
%   3. \@maketitle temporarily sets \thefootnote to \@fnsymbol; at replay time
%      that redefinition is gone, so \footnotetext[N] would use arabic numbering
%      and the marks in the title and the text at the bottom would not match.
% Fix: store the mark string and affiliation text in globally-defined csname
% macros indexed by a counter, then replay with \@footnotetext directly.
\makeatletter
\newcount\ts@affilcnt
\newcount\ts@loopcnt
\renewcommand{\thanks}[1]{%
  \stepcounter{footnote}%
  \protected@xdef\@thefnmark{\thefootnote}%
  \@footnotemark
  \global\advance\ts@affilcnt\@ne
  % xdef: globally store the current mark string (e.g. * or †) as a literal
  \expandafter\xdef\csname ts@affilM\the\ts@affilcnt\endcsname{\@thefnmark}%
  % gdef: globally store the affiliation text verbatim (no expansion)
  \expandafter\gdef\csname ts@affilT\the\ts@affilcnt\endcsname{#1}%
}
\newcommand\tsreplayfn{%
  \ts@loopcnt\z@
  \loop
    \advance\ts@loopcnt\@ne
    \ifnum\ts@loopcnt>\ts@affilcnt\else
      \xdef\@thefnmark{\csname ts@affilM\the\ts@loopcnt\endcsname}%
      \@footnotetext{\csname ts@affilT\the\ts@loopcnt\endcsname}%
  \repeat
}
\makeatother

\begin{document}

% Two-column layout: \maketitle is called outside the twocolumn class option so
% it does not internally invoke \twocolumn — \@thanks remains populated and is
% flushed explicitly after the block so affiliation footnotes are preserved.
\twocolumn[{%
\maketitle
}]
\tsreplayfn

\thispagestyle{plain}
\begin{callout}[callout note]{Note}
Highlights extra information that's useful but not critical, helping readers understand nuances.
\end{callout}
\begin{callout}[callout tip]{Tip}
Offers a practical hint that makes the task easier, saving the reader time or effort.
\end{callout}
\begin{callout}[callout warning]{Warning}
Draws attention to something that could cause problems if ignored, helping avoid mistakes.
\end{callout}
\begin{callout}[callout caution]{Caution}
Signals a potentially risky action, encouraging the reader to proceed carefully.
\end{callout}
\begin{callout}[callout important]{Important}
Emphasizes key information the reader must not overlook to ensure proper understanding.
\end{callout}
\begin{callout}[callout danger]{Danger}
Flags a serious hazard that could break things or cause real harm if mishandled.
\end{callout}
\begin{callout}[callout info]{Info}
Provides neutral, factual context that supports understanding without urgency.
\end{callout}
\begin{callout}[callout hint]{Hint}
Gives a gentle clue that helps the reader figure something out without revealing everything.
\end{callout}
\begin{callout}[callout info]{Seealso}
Points to related material so the reader can explore deeper or connected topics.
\end{callout}
\begin{callout}[callout question]{Question}
Raises an inquiry that prompts reflection or introduces a point the reader should consider.
\end{callout}
\begin{callout}[callout abstract]{Abstract}
Summarizes the core ideas to help the reader grasp the purpose of a section or document quickly.
\end{callout}
\begin{callout}[callout summary]{Summary}
Recaps key points so the reader can retain the most important information at a glance.
\end{callout}
\begin{callout}[callout success]{Success}
Celebrates a positive outcome or achievement, reinforcing good practices and results.
\end{callout}
\begin{callout}[callout failure]{Failure}
Highlights a setback or error, helping the reader learn from mistakes and avoid them in the future
\end{callout}
\begin{callout}[callout bug]{Bug}
Identifies a known issue or problem, guiding the reader on what to watch out for.
\end{callout}
\begin{callout}[callout quote]{Quote}
Presents a relevant quotation that adds authority or perspective to the content.
\end{callout}
\section*{Custom Admonition}\label{custom-admonition}

\begin{callout}[callout unicorn]{Unicorn}
A custom admonition with a unicorn theme, adding a whimsical touch to the information presented.
\end{callout}

\end{document}
